%% AC01
\begin{frame}[t]{Du modèle local à la forme globale}
\coursanalyse
Au premier ordre, on approche une petite variation par une application linéaire.
\begin{cours}{Trois nouvelles questions}
\textbf{Contrôler :} quelle amplitude de variation peut-on garantir ?

\medskip
\textbf{Affiner :} que gagne-t-on avec le second ordre ?

\medskip
\textbf{Comparer :} comment reconnaître une fonction convexe ?
\end{cours}
\vfill Régularité, Taylor et convexité relient ces trois questions.
\note{Cours original, pages PDF 3,31,59,85,109. Résultats admis ; interprétation sans démonstration.}
\end{frame}

%% AC02
\begin{frame}[t]{Les objets du premier ordre}
\coursanalyse
\begin{cours}{Une notation, un rôle}
\textbf{Différentielle :} $\df f{\vec a}$ est une application linéaire.

\medskip
\textbf{Jacobienne :} $J_f(\vec a)$ est sa matrice dans les bases canoniques.

\medskip
\textbf{Gradient, pour $f$ réelle :} $\grad f(\vec a)$ est un vecteur colonne.
\end{cours}
\[\df f{\vec a}(\vec h)=\ps{\grad f(\vec a)}{\vec h}.\]
\vfill Dans ce chapitre, les espaces sont euclidiens canoniques réels.
\note{Cours original, pages PDF 3,31. Résultats admis ; interprétation sans démonstration.}
\end{frame}

%% AC03
\begin{frame}[t]{Suivre un segment pour étudier une variation}
\coursanalyse
Si $[\vec a,\vec b]\subset U$, on pose
\[\vec\gamma(t)=\vec a+t(\vec b-\vec a),\qquad\varphi=f\circ\vec\gamma.\]
\begin{theoreme}{Si $f$ est différentiable sur l'ouvert $U$}
\[\varphi'(t)=\df f{\vec\gamma(t)}(\vec b-\vec a),\qquad 0<t<1.\]
\end{theoreme}
\begin{center}\begin{tikzpicture}
\fill(0,0)circle(2pt)node[below]{$\vec a$};\fill(7,0)circle(2pt)node[below]{$\vec b$};\draw[IMTCyan,thick](0,0)--(7,0);\fill[IMTOrange](3,0)circle(2pt)node[above]{$\vec\gamma(t)$};\draw[-{Stealth},IMTNavy](1,-.8)--(5,-.8)node[midway,below]{$\vec b-\vec a$};
\end{tikzpicture}\end{center}
\note{Cours original, pages PDF 3–4. Résultats admis ; interprétation sans démonstration.}
\end{frame}

%% AC04
\begin{frame}[t]{Un ensemble convexe contient ses segments}
\coursanalyse
\begin{notion}{Ensemble convexe}
$C\subset\Rn n$ est convexe si, pour tous $\vec a,\vec b\in C$,
\[\forall t\in[0,1],\qquad(1-t)\vec a+t\vec b\in C.\]
\end{notion}
\begin{center}\begin{tikzpicture}[scale=.7]
\fill[IMTCyan!12](0,0)ellipse[x radius=2,y radius=1.1];\draw[IMTCyan](0,0)ellipse[x radius=2,y radius=1.1];\draw[IMTNavy,thick](-1,-.3)--(1,.5);\fill(-1,-.3)circle(2pt);\fill(1,.5)circle(2pt);\node at(0,-1.6){segment inclus};
\draw[IMTCyan,thick](7,0)circle(1.1);\draw[IMTOrange,thick](5.9,0)--(8.1,0);\fill(5.9,0)circle(2pt);\fill(8.1,0)circle(2pt);\node at(7,-1.6){un cercle n'est pas convexe};
\end{tikzpicture}\end{center}
\vfill Les boules et les rectangles sont convexes. Un ouvert n'est pas forcément convexe.
\note{Cours original, pages PDF 13,109–112. Résultats admis ; interprétation sans démonstration.}
\end{frame}

%% AC05
\begin{frame}[t]{La norme d’une application linéaire}
\coursanalyse
\begin{notion}{Norme opérateur}
Pour $\Lin L:\Rn n\to\Rn p$ linéaire,
\[\norm{\Lin L}=\sup_{\norm{\vec h}=1}\norm{\Lin L(\vec h)}.\]
\end{notion}
\begin{theoreme}{Contrôler une image}
\[\forall\vec h,\qquad\norm{\Lin L(\vec h)}\le\norm{\Lin L}\,\norm{\vec h}.\]
\end{theoreme}
\vfill La norme opérateur mesure le plus grand facteur d'amplification des longueurs.
\note{Cours original, pages PDF 5–9. Résultats admis ; interprétation sans démonstration.}
\end{frame}

%% AC06
\begin{frame}[t]{Voir l’amplification maximale}
\coursanalyse
\begin{exemple}{Une application diagonale}
\[\Lin L\!\left(\vcol{x\\y}\right)=\vcol{-5x\\2y}.\]
La norme de l'image vérifie $\norm{\Lin L(\vec h)}\le5\norm{\vec h}$.
\end{exemple}
\begin{resultat}{Une direction atteint la borne}
$\vec e_1=\vcol{1\\0}$ est unitaire et $\norm{\Lin L(\vec e_1)}=5$.

\medskip
Donc $\norm{\Lin L}=5$.
\end{resultat}
\note{Cours original, pages PDF 5–9. Résultats admis ; interprétation sans démonstration.}
\end{frame}

%% AC07
\begin{frame}[t]{Une borne à partir des coefficients}
\coursanalyse
\begin{theoreme}{Matrice $A=(a_{ij})$ dans les bases canoniques}
La norme opérateur euclidienne satisfait
\[\norm A\le\sqrt{\sum_{i,j}a_{ij}^2}.\]
\end{theoreme}
\begin{cours}{Comment utiliser la borne ?}
Elle donne rapidement une majoration sans chercher la direction d'amplification maximale.
\end{cours}
\vfill Le membre de droite n'est pas en général la valeur exacte de la norme opérateur.
\note{Cours original, pages PDF 9–11. Résultats admis ; interprétation sans démonstration.}
\end{frame}

%% AC08
\begin{frame}[t]{Accroissements finis : une égalité scalaire}
\coursanalyse
\begin{theoreme}{Fonction à valeurs réelles}
Soient $f:U\to\R$ différentiable sur un ouvert $U$, $\vec a\ne\vec b$ et $[\vec a,\vec b]\subset U$.
Il existe $\theta\in]0,1[$ tel que, pour $\vec c=\vec a+\theta(\vec b-\vec a)$,
\[f(\vec b)-f(\vec a)=\df f{\vec c}(\vec b-\vec a).\]
\end{theoreme}
\begin{resultat}{Avec le gradient}
\[f(\vec b)-f(\vec a)=\ps{\grad f(\vec c)}{\vec b-\vec a}.\]
\end{resultat}
\note{Cours original, pages PDF 4–5. Résultats admis ; interprétation sans démonstration.}
\end{frame}

%% AC09
\begin{frame}[t]{Accroissements finis : une borne vectorielle}
\coursanalyse
\begin{theoreme}{Inégalité des accroissements finis}
Soient $f:U\to\Rn p$ différentiable sur l'ouvert $U$ et $[\vec a,\vec b]\subset U$.
Si $\norm{\df f{\vec x}}\le M$ sur ce segment, alors
\[\norm{f(\vec b)-f(\vec a)}\le M\norm{\vec b-\vec a}.\]
\end{theoreme}
\begin{resultat}{Cas réel}
On peut remplacer la borne sur la différentielle par $\norm{\grad f(\vec x)}\le M$.
\end{resultat}
\vfill Pour une sortie vectorielle, ne pas supposer une égalité en un unique point intermédiaire.
\note{Cours original, pages PDF 11–13. Résultats admis ; interprétation sans démonstration.}
\end{frame}

%% AC10
\begin{frame}[t]{Être lipschitzienne : une variation contrôlée}
\coursanalyse
\begin{notion}{Constante $M\ge0$}
$f:C\to\Rn p$ est $M$-lipschitzienne si
\[\forall\vec x,\vec y\in C,\quad\norm{f(\vec y)-f(\vec x)}\le M\norm{\vec y-\vec x}.\]
\end{notion}
\begin{resultat}{Une garantie quantitative}
Une variation d'entrée de taille $\delta$ provoque une variation de sortie d'au plus $M\delta$.
\end{resultat}
\vfill Toute fonction lipschitzienne est continue ; la réciproque est fausse.
\note{Cours original, pages PDF 13–14. Résultats admis ; interprétation sans démonstration.}
\end{frame}

%% AC11
\begin{frame}[t]{Une borne locale devient un contrôle global}
\coursanalyse
\begin{theoreme}{Sur un ensemble convexe $K$}
Soit $f$ différentiable sur un ouvert contenant $K$.
Si $\norm{\df f{\vec x}}\le M$ pour tout $\vec x\in K$, alors $f$ est $M$-lipschitzienne sur $K$.
\end{theoreme}
\begin{cours}{Les deux vérifications}
\textbf{Géométrie :} les segments entre points de $K$ restent dans $K$.

\medskip
\textbf{Calcul :} la même constante $M$ majore la différentielle partout sur $K$.
\end{cours}
\vfill $K$ peut être une boule fermée ou un rectangle fermé : l'ouvert porte sur le domaine où $f$ est différentiable.
\note{Cours original, pages PDF 13–15. Résultats admis ; interprétation sans démonstration.}
\end{frame}

%% AC12
\begin{frame}[t]{Exemple : contrôler une variation sur une boule}
\coursanalyse
\begin{exemple}{$f(x,y)=x^2+4y^2$, sur $K=\overline B(\zero,2)$}
\[\grad f(x,y)=\vcol{2x\\8y},\qquad\norm{\grad f(x,y)}\le8\norm{\vec x}\le16.\]
La boule $K$ est convexe et $f$ est différentiable sur $\Rn2$.
\end{exemple}
\begin{resultat}{Une borne valable pour tous les couples}
\[|f(\vec b)-f(\vec a)|\le16\norm{\vec b-\vec a},\qquad\vec a,\vec b\in K.\]
\end{resultat}
\note{Cours original, pages PDF 11–15. Résultats admis ; interprétation sans démonstration.}
\end{frame}

%% AC13
\begin{frame}[t]{Une différentielle nulle force la constance}
\coursanalyse
\begin{theoreme}{Sur un ouvert convexe non vide $U$}
Si $f:U\to\Rn p$ est différentiable et
\[\forall\vec x\in U,\quad\df f{\vec x}=\Lin0,\]
alors $f$ est constante sur $U$.
\end{theoreme}
\begin{resultat}{Pour une fonction réelle}
Un gradient nul \textbf{partout sur $U$} implique la constance.
\end{resultat}
\vfill Un gradient nul en un seul point ne suffit pas.
\note{Cours original, pages PDF 14–15. Résultats admis ; interprétation sans démonstration.}
\end{frame}

%% AC14
\begin{frame}[t]{La classe C1 : une différentielle continue}
\coursanalyse
\begin{notion}{Classe $\mathcal C^1$ sur un ouvert $U$}
$f$ est de classe $\mathcal C^1$ si elle est différentiable sur $U$ et si
\[\vec x\longmapsto\df f{\vec x}\]
est continue pour la norme opérateur.
\end{notion}
\begin{resultat}{Le sens de la régularité}
Le modèle linéaire varie continûment lorsque le point de départ se déplace.
\end{resultat}
\vfill Classe $\mathcal C^1$ implique différentiabilité ; la réciproque est fausse.
\note{Cours original, pages PDF 31–33. Résultats admis ; interprétation sans démonstration.}
\end{frame}

%% AC15
\begin{frame}[t]{Le critère pratique de classe C1}
\coursanalyse
\begin{theoreme}{Dérivées partielles continues sur un ouvert}
$f:U\to\Rn p$ est de classe $\mathcal C^1$ si et seulement si toutes les dérivées partielles de ses composantes existent et sont continues sur $U$.
\end{theoreme}
\begin{cours}{Une méthode en trois étapes}
\textbf{1.} Calculer chaque dérivée partielle.

\medskip
\textbf{2.} Justifier sa continuité sur le domaine.

\medskip
\textbf{3.} Conclure $\mathcal C^1$, donc différentiable.
\end{cours}
\note{Cours original, pages PDF 33–39. Résultats admis ; interprétation sans démonstration.}
\end{frame}

%% AC16
\begin{frame}[t]{Un exemple de régularité}
\coursanalyse
\begin{exemple}{$F(x,y)=\vcol{xe^y\\x^2+\cos y}$}
\[J_F(x,y)=\begin{pmatrix}e^y&xe^y\\2x&-\sin y\end{pmatrix}.\]
Les quatre coefficients sont des fonctions continues sur $\Rn2$.
\end{exemple}
\begin{resultat}{La conclusion utile}
$F$ est de classe $\mathcal C^1$ sur $\Rn2$ ; on peut donc utiliser cette jacobienne pour calculer sa différentielle.
\end{resultat}
\note{Cours original, pages PDF 34–43. Résultats admis ; interprétation sans démonstration.}
\end{frame}

%% AC17
\begin{frame}[t]{Ce que l’existence des partielles ne dit pas}
\coursanalyse
\begin{cours}{Trois niveaux à distinguer}
\textbf{Partielles existantes :} les variations selon les axes existent.

\medskip
\textbf{Différentiabilité :} un même modèle linéaire contrôle toutes les directions.

\medskip
\textbf{Classe $\mathcal C^1$ :} ce modèle varie continûment avec le point.
\end{cours}
\begin{resultat}{Une condition suffisante pratique}
Des partielles continues sur un ouvert donnent la classe $\mathcal C^1$, et donc la différentiabilité.
\end{resultat}
\note{Cours original, pages PDF 34–39. Résultats admis ; interprétation sans démonstration.}
\end{frame}

%% AC18
\begin{frame}[t]{Construire des fonctions de classe C1}
\coursanalyse
\begin{theoreme}{Règles de stabilité}
\textbf{Somme et produit :} ils conservent la classe $\mathcal C^1$ pour des fonctions réelles.

\medskip
\textbf{Quotient :} sur le domaine où le dénominateur ne s'annule pas.

\medskip
\textbf{Composition :} $g\circ f$ est de classe $\mathcal C^1$ si $f$ et $g$ le sont sur des ouverts compatibles.
\end{theoreme}
\vfill Les formules du premier ordre restent les mêmes ; on contrôle en plus leur continuité.
\note{Cours original, pages PDF 40–43. Résultats admis ; interprétation sans démonstration.}
\end{frame}

%% AC19
\begin{frame}[t]{Pourquoi regarder le second ordre ?}
\coursanalyse
\begin{cours}{Après l’approximation affine}
La différentielle première décrit la pente locale.

\medskip
La différentielle seconde décrit comment cette pente varie.
\end{cours}
\begin{center}\begin{tikzpicture}[x=1.25cm,y=.65cm]
\draw[axes](-2.2,0)--(2.2,0);\draw[axes](0,-.3)--(0,3.7);
\draw[courb,domain=-1.8:1.8,samples=50]plot(\x,{\x*\x});\draw[IMTOrange,thick](-1.8,0)--(1.8,0);
\node[anchor=west,align=left]at(3,2){$f(t)=t^2$\\[2mm]$f'(0)=0$, mais $f''(0)=2$.};
\end{tikzpicture}\end{center}
\vfill Le second ordre apporte une information absente de la tangente.
\note{Cours original, pages PDF 59–60,85–90. Résultats admis ; interprétation sans démonstration.}
\end{frame}

%% AC20
\begin{frame}[t]{La différentielle seconde agit sur deux directions}
\coursanalyse
\begin{notion}{Deux fois différentiable en $\vec a$}
L'application $\vec x\mapsto\df f{\vec x}$ est différentiable en $\vec a$ ; $f$ est différentiable dans un voisinage de ce point.
\end{notion}
\begin{notion}{Différentielle seconde}
\[\ds f{\vec a}(\vec h,\vec k)
=\bigl(\df{(df)}{\vec a}(\vec h)\bigr)(\vec k).\]
Elle est \textbf{bilinéaire} : linéaire dans chacun des deux déplacements séparément.
\end{notion}
\vfill Le premier déplacement fait varier la différentielle ; le second évalue cette variation.
\note{Cours original, pages PDF 59–61. Résultats admis ; interprétation sans démonstration.}
\end{frame}

%% AC21
\begin{frame}[t]{L’ordre des dérivations partielles}
\coursanalyse
\begin{notion}{Dérivée partielle seconde}
\[\frac{\partial^2 f}{\partial x_i\partial x_j}
=\frac{\partial}{\partial x_i}\left(\frac{\partial f}{\partial x_j}\right).\]
On dérive \textbf{d'abord selon $x_j$}, puis selon $x_i$.
\end{notion}
\begin{exemple}{Pour $f(x,y)=x^2y$}
\[\frac{\partial f}{\partial y}=x^2,\qquad
\frac{\partial^2f}{\partial x\partial y}=2x.\]
\end{exemple}
\vfill On n'échange pas encore les deux dérivations sans hypothèse de régularité.
\note{Cours original, pages PDF 61–62. Résultats admis ; interprétation sans démonstration.}
\end{frame}

%% AC22
\begin{frame}[t]{La Hessienne rassemble les dérivées secondes}
\coursanalyse
\begin{notion}{Fonction à valeurs réelles}
\[H_f(\vec a)=\left(\frac{\partial^2f}{\partial x_i\partial x_j}(\vec a)\right)_{1\le i,j\le n}.\]
\end{notion}
En deux variables :
\[H_f=\begin{pmatrix}
\dfrac{\partial^2 f}{\partial x^2}&\dfrac{\partial^2 f}{\partial x\partial y}\\
\dfrac{\partial^2 f}{\partial y\partial x}&\dfrac{\partial^2 f}{\partial y^2}
\end{pmatrix}.\]
\vfill Une Hessienne est une matrice carrée de taille $n$ ; sa symétrie sera assurée par le théorème de Schwarz.
\note{Cours original, pages PDF 61–64. Résultats admis ; interprétation sans démonstration.}
\end{frame}

%% AC23
\begin{frame}[t]{La classe C2 : un critère calculable}
\coursanalyse
\begin{notion}{Classe $\mathcal C^2$ sur un ouvert $U$}
$f$ est différentiable et l'application $\vec x\mapsto\df f{\vec x}$ est de classe $\mathcal C^1$.
\end{notion}
\begin{theoreme}{Critère par les composantes}
$f:U\to\Rn p$ est de classe $\mathcal C^2$ si et seulement si toutes les dérivées partielles secondes de ses composantes existent et sont continues sur $U$.
\end{theoreme}
\vfill Une fonction de classe $\mathcal C^2$ est aussi de classe $\mathcal C^1$.
\note{Cours original, pages PDF 64–66. Résultats admis ; interprétation sans démonstration.}
\end{frame}

%% AC24
\begin{frame}[t]{Schwarz : quand les dérivations s’échangent}
\coursanalyse
\begin{theoreme}{Sous l’hypothèse $\mathcal C^2$}
Si $f$ est de classe $\mathcal C^2$ sur un ouvert $U$,
\[\frac{\partial^2f}{\partial x_i\partial x_j}(\vec a)
=\frac{\partial^2f}{\partial x_j\partial x_i}(\vec a),\qquad\vec a\in U.\]
\end{theoreme}
\begin{resultat}{Pour une fonction réelle}
La Hessienne est \textbf{symétrique} : $H_f(\vec a)^\top=H_f(\vec a)$.
\end{resultat}
\vfill La régularité justifie l'échange ; une simple ressemblance des formules ne remplace pas cette hypothèse.
\note{Cours original, pages PDF 66–68. Résultats admis ; interprétation sans démonstration.}
\end{frame}

%% AC25
\begin{frame}[t]{Un exemple de Hessienne}
\coursanalyse
\begin{exemple}{$f(x,y)=x^2+xy+2y^2+y^3$}
\[\grad f(x,y)=\vcol{2x+y\\x+4y+3y^2},\qquad
H_f(x,y)=\begin{pmatrix}2&1\\1&4+6y\end{pmatrix}.\]
\end{exemple}
\begin{resultat}{Au point $\zero$}
\[H_f(\zero)=\begin{pmatrix}2&1\\1&4\end{pmatrix}.\]
\end{resultat}
\vfill La fonction est polynomiale ; ses dérivées secondes sont continues.
\note{Cours original, pages PDF 61–68. Résultats admis ; interprétation sans démonstration.}
\end{frame}

%% AC26
\begin{frame}[t]{Relier différentielle seconde et Hessienne}
\coursanalyse
\begin{theoreme}{Fonction réelle de classe $\mathcal C^2$}
Dans les bases canoniques,
\[\ds f{\vec a}(\vec h,\vec k)=\vec h^{\top}H_f(\vec a)\vec k.\]
\end{theoreme}
\begin{exemple}{Pour $H_f(\zero)=\left(\begin{smallmatrix}2&1\\1&4\end{smallmatrix}\right)$}
Avec $\vec h=\vcol{h_1\\h_2}$ et $\vec k=\vcol{k_1\\k_2}$,
\[\ds f{\zero}(\vec h,\vec k)=2h_1k_1+h_1k_2+h_2k_1+4h_2k_2.\]
\end{exemple}
\note{Cours original, pages PDF 62–64. Résultats admis ; interprétation sans démonstration.}
\end{frame}

%% AC27
\begin{frame}[t]{Une même direction fait apparaître un carré}
\coursanalyse
\begin{notion}{Forme quadratique associée}
\[\ds f{\vec a}(\vec h,\vec h)=\vec h^{\top}H_f(\vec a)\vec h.\]
\end{notion}
\begin{exemple}{Dans l’exemple précédent, au point $\zero$}
\[\vec h^{\top}H_f(\zero)\vec h
=2h_1^2+2h_1h_2+4h_2^2.\]
\end{exemple}
\vfill Le coefficient du terme croisé reçoit les \textbf{deux} contributions hors diagonale.
\note{Cours original, pages PDF 62–64. Résultats admis ; interprétation sans démonstration.}
\end{frame}

%% AC28
\begin{frame}[t]{Taylor–Young au second ordre}
\coursanalyse
\begin{theoreme}{Si $f:U\to\R$ est de classe $\mathcal C^2$}
Pour $\vec a\in U$ ouvert et $\vec h\to\zero$,
\[\begin{aligned}
f(\vec a+\vec h)={}&f(\vec a)+\df f{\vec a}(\vec h)\\
&+\frac12\ds f{\vec a}(\vec h,\vec h)+o(\norm{\vec h}^2).
\end{aligned}\]
\end{theoreme}
\begin{resultat}{Ce qui s’ajoute au premier ordre}
Un terme quadratique, avec le facteur $1/2$ ; le reste est négligeable devant le carré de la norme.
\end{resultat}
\note{Cours original, pages PDF 88–90. Résultats admis ; interprétation sans démonstration.}
\end{frame}

%% AC29
\begin{frame}[t]{Lire le modèle quadratique}
\coursanalyse
\begin{resultat}{Version gradient et Hessienne}
\[f(\vec a+\vec h)\simeq f(\vec a)+\ps{\grad f(\vec a)}{\vec h}
+\frac12\vec h^\top H_f(\vec a)\vec h.\]
\end{resultat}
\begin{cours}{Trois informations au même point}
\textbf{Valeur :} $f(\vec a)$.

\medskip
\textbf{Variation linéaire :} le gradient en $\vec a$.

\medskip
\textbf{Variation quadratique :} la Hessienne en $\vec a$.
\end{cours}
\note{Cours original, pages PDF 88–90. Résultats admis ; interprétation sans démonstration.}
\end{frame}

%% AC30
\begin{frame}[t]{Un exemple de Taylor au second ordre}
\coursanalyse
\begin{exemple}{$f(x,y)=x^2+xy+2y^2+y^3$, au point $\zero$}
$f(\zero)=0$, $\grad f(\zero)=\zero$ et $H_f(\zero)=\left(\begin{smallmatrix}2&1\\1&4\end{smallmatrix}\right)$.
\[f(h,k)=\underbrace{h^2+hk+2k^2}_{\text{modèle au second ordre}}+k^3.\]
Pour $\vec h=\vcol{0{,}02\\-0{,}01}$, le modèle vaut $0{,}0004$.
\end{exemple}
\begin{resultat}{Interpréter l’écart}
La valeur exacte est $0{,}000399$ : le reste cubique vaut $-0{,}000001$.
\end{resultat}
\note{Cours original, pages PDF 88–90. Résultats admis ; interprétation sans démonstration.}
\end{frame}

%% AC31
\begin{frame}[t]{Taylor avec un point intermédiaire}
\coursanalyse
\begin{theoreme}{Version exacte, fonction réelle de classe $\mathcal C^2$}
Si $[\vec a,\vec a+\vec h]\subset U$, il existe $\theta\in]0,1[$ tel que
\[\begin{aligned}
f(\vec a+\vec h)={}&f(\vec a)+\df f{\vec a}(\vec h)\\
&+\frac12\vec h^\top H_f(\vec a+\theta\vec h)\vec h.
\end{aligned}\]
\end{theoreme}
\begin{resultat}{Ne pas mélanger les deux formules}
Taylor–Young utilise la Hessienne en $\vec a$ avec un reste.
La formule exacte utilise une Hessienne en un point intermédiaire.
\end{resultat}
\note{Cours original, pages PDF 87–88. Résultats admis ; interprétation sans démonstration.}
\end{frame}

%% AC32
\begin{frame}[t]{Contrôler l’erreur de linéarisation}
\coursanalyse
\begin{theoreme}{Hessienne bornée sur le segment}
Soit $f:U\to\R$ de classe $\mathcal C^2$, avec $[\vec a,\vec a+\vec h]\subset U$.
Si $\norm{H_f(\vec x)}\le M$ sur ce segment, alors
\[\left|f(\vec a+\vec h)-f(\vec a)-\df f{\vec a}(\vec h)\right|
\le\frac M2\norm{\vec h}^2.\]
\end{theoreme}
\vfill Une Hessienne uniformément bornée donne une erreur \textbf{quadratique} pour l'approximation affine.
\note{Cours original, pages PDF 93–94. Résultats admis ; interprétation sans démonstration.}
\end{frame}

%% AC33
\begin{frame}[t]{Définir une fonction convexe}
\coursanalyse
\begin{notion}{Sur un ensemble convexe $C$}
$f:C\to\R$ est convexe si, pour tous $\vec x,\vec y\in C$ et $t\in[0,1]$,
\[f((1-t)\vec x+t\vec y)\le(1-t)f(\vec x)+tf(\vec y).\]
\end{notion}
\begin{resultat}{Le sens de l’inégalité}
La valeur au point intermédiaire ne dépasse pas la moyenne pondérée des valeurs aux deux extrémités.
\end{resultat}
\vfill Le domaine doit lui-même contenir les segments utilisés.
\note{Cours original, pages PDF 113–114. Résultats admis ; interprétation sans démonstration.}
\end{frame}

%% AC34
\begin{frame}[t]{La courbe reste sous la corde}
\coursanalyse
\begin{center}\begin{tikzpicture}[x=1.3cm,y=.7cm]
\draw[axes](-.3,0)--(3.4,0)node[right]{$x$};\draw[axes](0,-.3)--(0,5.2);
\draw[courb,domain=0:3,samples=60]plot(\x,{.5*\x*\x});
\draw[IMTOrange,thick](.5,.125)--(2.5,3.125);
\fill(.5,.125)circle(2pt)node[below]{$a$};\fill(2.5,3.125)circle(2pt)node[above left]{$b$};
\draw[pointille](1.5,1.125)--(1.5,1.625);\fill[IMTCyan](1.5,1.125)circle(2pt);\fill[IMTOrange](1.5,1.625)circle(2pt);
\node[anchor=west,align=left]at(4,3){\textcolor{IMTOrange}{corde : moyenne des valeurs}\\[3mm]\textcolor{IMTNavy}{courbe : valeur intermédiaire}};
\end{tikzpicture}\end{center}
\begin{resultat}{Pour plusieurs variables}
La même comparaison se fait sur chaque segment du domaine.
\end{resultat}
\note{Cours original, pages PDF 113–114. Résultats admis ; interprétation sans démonstration.}
\end{frame}

%% AC35
\begin{frame}[t]{Des exemples et un contre-exemple}
\coursanalyse
\begin{exemple}{Fonctions convexes}
\textbf{Affine :} $f(\vec x)=\ps{\vec c}{\vec x}+d$ ; l'inégalité est une égalité.

\medskip
\textbf{Quadratique :} $f(\vec x)=\norm{\vec x}^2$.
\end{exemple}
\begin{exemple}{Une fonction qui n’est pas convexe}
$f(t)=-t^2$ : aux points $-1$ et $1$, la moyenne des valeurs vaut $-1$, alors que $f(0)=0$.
\end{exemple}
\vfill La convexité concerne toutes les paires et tous les poids, pas quelques points seulement.
\note{Cours original, pages PDF 113–115. Résultats admis ; interprétation sans démonstration.}
\end{frame}

%% AC36
\begin{frame}[t]{Le critère du premier ordre : un support affine}
\coursanalyse
\begin{theoreme}{Fonction différentiable sur un ouvert convexe $U$}
$f:U\to\R$ est convexe si et seulement si, pour tous $\vec x,\vec y\in U$,
\[f(\vec y)\ge f(\vec x)+\ps{\grad f(\vec x)}{\vec y-\vec x}.\]
\end{theoreme}
\begin{resultat}{Lecture géométrique}
L'approximation affine tangente est un \textbf{minorant global} de la fonction sur $U$.
\end{resultat}
\vfill Pour une fonction convexe : sous les cordes, au-dessus des tangentes.
\note{Cours original, pages PDF 121–123. Résultats admis ; interprétation sans démonstration.}
\end{frame}

%% AC37
\begin{frame}[t]{Le critère de convexité par la Hessienne}
\coursanalyse
\begin{theoreme}{Fonction réelle de classe $\mathcal C^2$}
Sur un \textbf{ouvert convexe} $U$,
\[f\text{ convexe}\quad\Longleftrightarrow\quad
\forall\vec x\in U,\ \forall\vec h,\quad\vec h^\top H_f(\vec x)\vec h\ge0.\]
\end{theoreme}
\begin{resultat}{Le lien avec l’algèbre}
Il faut que la Hessienne soit \textbf{positive semi-définie en tout point} du domaine.
\end{resultat}
\vfill Une valeur négative dans une direction et en un point suffit à exclure la convexité sur $U$.
\note{Cours original, pages PDF 124–125. Résultats admis ; interprétation sans démonstration.}
\end{frame}

%% AC38
\begin{frame}[t]{Lire le signe d’une matrice symétrique}
\coursanalyse
\begin{cours}{Pour une matrice réelle symétrique $H$}
\textbf{Positive semi-définie :} $\vec h^\top H\vec h\ge0$ pour tout $\vec h$.

\medskip
\textbf{Définie positive :} $\vec h^\top H\vec h>0$ pour tout $\vec h\ne\zero$.
\end{cours}
\begin{theoreme}{Critère spectral}
Semi-définie positive : toutes les valeurs propres sont $\ge0$.

\medskip
Définie positive : toutes les valeurs propres sont $>0$.
\end{theoreme}
\vfill Une décomposition en carrés peut éviter le calcul des valeurs propres.
\note{Cours original, pages PDF 124–126 ; prérequis d’algèbre. Résultats admis ; interprétation sans démonstration.}
\end{frame}

%% AC39
\begin{frame}[t]{Exemple : la convexité d’une quadratique}
\coursanalyse
\begin{exemple}{$q(x,y)=x^2+xy+2y^2+3x$}
\[H_q=\begin{pmatrix}2&1\\1&4\end{pmatrix}.\]
Pour $\vec h=\vcol{h\\k}$,
\[\vec h^\top H_q\vec h=2\left(h+\frac k2\right)^2+\frac72k^2\ge0.\]
\end{exemple}
\begin{resultat}{Conclusion sur tout le domaine}
La Hessienne est positive semi-définie partout ; $q$ est convexe sur $\Rn2$.
\end{resultat}
\note{Cours original, pages PDF 124–125. Résultats admis ; interprétation sans démonstration.}
\end{frame}

%% AC40
\begin{frame}[t]{La stricte convexité}
\coursanalyse
\begin{notion}{Une inégalité stricte entre points distincts}
Sur un ensemble convexe $C$, $f$ est strictement convexe si, pour $\vec x\ne\vec y$ et $0<t<1$,
\[f((1-t)\vec x+t\vec y)<(1-t)f(\vec x)+tf(\vec y).\]
\end{notion}
\begin{resultat}{Ce qui change}
Entre deux points distincts, la valeur intermédiaire est strictement sous la corde.
\end{resultat}
\vfill Les extrémités $t=0$ et $t=1$ sont exclues de l'inégalité stricte.
\note{Cours original, pages PDF 114,126–127. Résultats admis ; interprétation sans démonstration.}
\end{frame}

%% AC41
\begin{frame}[t]{Hessienne définie positive : un critère suffisant}
\coursanalyse
\begin{theoreme}{Sur un ouvert convexe $U$}
Si $f:U\to\R$ est de classe $\mathcal C^2$ et si $H_f(\vec x)$ est définie positive pour tout $\vec x\in U$, alors $f$ est strictement convexe.
\end{theoreme}
\begin{exemple}{La réciproque est fausse}
$f(t)=t^4$ est strictement convexe sur $\R$, mais $f''(0)=0$.
\end{exemple}
\vfill Ne pas transformer cette condition suffisante en équivalence.
\note{Cours original, pages PDF 126–127. Résultats admis ; interprétation sans démonstration.}
\end{frame}

%% AC42
\begin{frame}[t]{Stricte ne signifie pas courbure uniforme}
\coursanalyse
\begin{exemple}{Deux fonctions strictement convexes}
\textbf{$f(t)=t^4$ :} $f''(t)=12t^2$ s'annule en $0$.

\medskip
\textbf{$g(t)=e^t$ :} $g''(t)=e^t>0$ partout, mais tend vers $0$ lorsque $t\to-\infty$.
\end{exemple}
\begin{resultat}{La question suivante}
Peut-on imposer un même niveau de courbure strictement positif sur tout le domaine ?
\end{resultat}
\note{Cours original, pages PDF 126–129. Résultats admis ; interprétation sans démonstration.}
\end{frame}

%% AC43
\begin{frame}[t]{La forte convexité : un écart quantifié}
\coursanalyse
\begin{notion}{Paramètre uniforme $\mu>0$}
Sur un ensemble convexe $C$, $f$ est $\mu$-fortement convexe si, pour tous $\vec x,\vec y\in C$ et $t\in[0,1]$,
\[\begin{aligned}
f((1-t)\vec x+t\vec y)\le{}&(1-t)f(\vec x)+tf(\vec y)\\
&-\frac\mu2t(1-t)\norm{\vec y-\vec x}^2.
\end{aligned}\]
\end{notion}
\vfill Le même $\mu$ doit convenir à tous les points et tous les segments du domaine.
\note{Cours original, pages PDF 127–128. Résultats admis ; interprétation sans démonstration.}
\end{frame}

%% AC44
\begin{frame}[t]{Forte convexité et Hessienne}
\coursanalyse
\begin{theoreme}{Fonction réelle $\mathcal C^2$ sur un ouvert convexe $U$}
Pour un paramètre fixé $\mu>0$, $f$ est $\mu$-fortement convexe si et seulement si
\[\forall\vec x\in U,\ \forall\vec h,\qquad
\vec h^\top H_f(\vec x)\vec h\ge\mu\norm{\vec h}^2.\]
\end{theoreme}
\begin{resultat}{Lecture matricielle}
$H_f(\vec x)-\mu I_n$ doit être positive semi-définie en tout point.
\end{resultat}
\vfill Des valeurs propres strictement positives point par point ne suffisent pas sans une même borne $\mu>0$.
\note{Cours original, pages PDF 128–129. Résultats admis ; interprétation sans démonstration.}
\end{frame}

%% AC45
\begin{frame}[t]{Trouver un paramètre de forte convexité}
\coursanalyse
\begin{exemple}{Pour $q(x,y)=x^2+xy+2y^2+3x$}
\[H_q-I_2=\begin{pmatrix}1&1\\1&3\end{pmatrix}.\]
Pour $\vec h=\vcol{h\\k}$,
\[\vec h^\top(H_q-I_2)\vec h=(h+k)^2+2k^2\ge0.\]
\end{exemple}
\begin{resultat}{Un paramètre valable}
$q$ est $1$-fortement convexe sur $\Rn2$.
Le paramètre $1$ est une borne commode ; il n'est pas annoncé comme optimal.
\end{resultat}
\note{Cours original, pages PDF 128–129. Résultats admis ; interprétation sans démonstration.}
\end{frame}

%% AC46
\begin{frame}[t]{Un minorant plus précis que la tangente}
\coursanalyse
\begin{theoreme}{Fonction différentiable sur un ouvert convexe}
Pour $\mu>0$, $f$ est $\mu$-fortement convexe si et seulement si, pour tous $\vec x,\vec y$ du domaine,
\[\begin{aligned}
f(\vec y)\ge{}&f(\vec x)+\ps{\grad f(\vec x)}{\vec y-\vec x}\\
&+\frac\mu2\norm{\vec y-\vec x}^2.
\end{aligned}\]
\end{theoreme}
\begin{resultat}{Une garantie quantitative}
La fonction reste au-dessus de son modèle affine, avec une marge quadratique uniforme.
\end{resultat}
\note{Cours original, pages PDF 128–129. Résultats admis ; interprétation sans démonstration.}
\end{frame}

%% AC47
\begin{frame}[t]{Convexe, stricte, forte : les implications}
\coursanalyse
\begin{center}\begin{tikzpicture}
\node[align=center]at(0,0){fortement\\convexe};\node[align=center]at(5,0){strictement\\convexe};\node at(10,0){convexe};
\draw[-{Stealth},thick,IMTCyan](1.5,0)--(3.4,0);\draw[-{Stealth},thick,IMTCyan](6.6,0)--(8.7,0);
\end{tikzpicture}\end{center}
\begin{cours}{Les réciproques sont fausses}
\textbf{Convexe sans être stricte :} une fonction affine sur $\R$.

\medskip
\textbf{Stricte sans être forte sur $\R$ :} $t\mapsto t^4$.
\end{cours}
\vfill La convexité forte dépend d'une constante uniforme et du domaine considéré.
\note{Cours original, pages PDF 114,126–129. Résultats admis ; interprétation sans démonstration.}
\end{frame}

%% AC48
\begin{frame}[t]{Une information en un point reste locale}
\coursanalyse
\begin{exemple}{$f(x,y)=x^2+xy+2y^2+y^3$}
$H_f(\zero)=\left(\begin{smallmatrix}2&1\\1&4\end{smallmatrix}\right)$ est définie positive.

\medskip
Mais au point $\vec b=\vcol{0\\-1}$,
\[H_f(\vec b)=\begin{pmatrix}2&1\\1&-2\end{pmatrix},\qquad
\vec e_2^\top H_f(\vec b)\vec e_2=-2.\]
\end{exemple}
\begin{resultat}{Conclusion}
La fonction n'est pas convexe sur $\Rn2$ : il faut contrôler la Hessienne \textbf{sur tout le domaine}.
\end{resultat}
\note{Cours original, pages PDF 124–129. Résultats admis ; interprétation sans démonstration.}
\end{frame}

%% AC49
\begin{frame}[t]{Une démarche d’étude en quatre questions}
\coursanalyse
\begin{cours}{Garder les hypothèses dans l’ordre}
\textbf{1. Domaine :} est-il ouvert ? Le domaine d'étude est-il convexe ?

\medskip
\textbf{2. Régularité :} les partielles nécessaires sont-elles continues ?

\medskip
\textbf{3. Objectif :} contrôler une variation, approcher par Taylor ou établir une convexité ?

\medskip
\textbf{4. Critère :} borne uniforme, développement au point ou signe de la Hessienne partout ?
\end{cours}
\vfill Ne pas remplacer une hypothèse globale par une vérification en un seul point.
\note{Cours original, pages PDF 3–129. Résultats admis ; interprétation sans démonstration.}
\end{frame}

%% AC50
\begin{frame}[t]{Les liens essentiels du chapitre}
\coursanalyse
\begin{cours}{Du local au global}
\textbf{Différentielle bornée + segments inclus :} contrôle des variations.

\medskip
\textbf{Partielles continues :} classe $\mathcal C^1$ et différentiabilité.

\medskip
\textbf{Second ordre continu :} Taylor et Hessienne symétrique.

\medskip
\textbf{Hessienne positive partout + domaine convexe :} convexité.

\medskip
\textbf{Minoration uniforme par $\mu I_n$ :} forte convexité.
\end{cours}
\vfill Les quatre pages suivantes sont des compléments séparés du socle.
\note{Cours original, pages PDF 3–129. Résultats admis ; interprétation sans démonstration.}
\end{frame}

%% AC51
\begin{frame}[t]{Complément — Les limites du TAF vectoriel}
\coursanalyse
\begin{exemple}{Une courbe fermée}
\[F(t)=\vcol{\cos t\\\sin t},\qquad t\in[0,2\pi].\]
$F(2\pi)-F(0)=\zero$, mais $F'(t)=\vcol{-\sin t\\\cos t}$ n'est jamais nul.
\end{exemple}
\begin{resultat}{Ce qui échoue}
Il n'existe aucun $c\in]0,2\pi[$ tel que
\[F(2\pi)-F(0)=2\pi F'(c).\]
\end{resultat}
\vfill Pour les sorties vectorielles, l'inégalité des accroissements finis reste l'outil général.
\note{Cours original, pages PDF 4–5,11–13. Résultats admis ; interprétation sans démonstration.}
\end{frame}

%% AC52
\begin{frame}[t]{Complément — Différentiable sans être C1}
\coursanalyse
\begin{exemple}{Une fonction réelle d’une variable}
\[f(t)=\begin{cases}t^2\sin(1/t),&t\ne0,\\0,&t=0.\end{cases}\]
$f$ est dérivable partout, avec $f'(0)=0$.

\medskip
Pour $t\ne0$, $f'(t)=2t\sin(1/t)-\cos(1/t)$.
\end{exemple}
\begin{resultat}{La limite manquante}
$f'$ oscille près de $0$ : elle n'est pas continue en $0$. Ainsi $f$ n'est pas de classe $\mathcal C^1$ au voisinage de $0$.
\end{resultat}
\note{Cours original, pages PDF 38–39. Résultats admis ; interprétation sans démonstration.}
\end{frame}

%% AC53
\begin{frame}[t]{Complément — Composer au second ordre}
\coursanalyse
\begin{theoreme}{Applications $\mathcal C^2$ sur des ouverts compatibles}
Posons $\vec u=\df f{\vec a}(\vec h)$ et $\vec v=\df f{\vec a}(\vec k)$.
\[\begin{aligned}
\ds{(g\circ f)}{\vec a}(\vec h,\vec k)
={}&\ds g{f(\vec a)}(\vec u,\vec v)\\
&+\df g{f(\vec a)}\bigl(\ds f{\vec a}(\vec h,\vec k)\bigr).
\end{aligned}\]
\end{theoreme}
\begin{resultat}{Deux contributions}
La variation propre de $g$, puis la variation du modèle de $f$ transmise à travers $g$.
\end{resultat}
\note{Cours original, pages PDF 69–71. Résultats admis ; interprétation sans démonstration.}
\end{frame}

%% AC54
\begin{frame}[t]{Complément — Jensen pour une moyenne pondérée}
\coursanalyse
\begin{theoreme}{Fonction convexe sur un ensemble convexe $C$}
Pour $\vec x_1,\ldots,\vec x_m\in C$ et des poids $\lambda_i\ge0$ de somme $1$,
\[f\!\left(\sum_{i=1}^m\lambda_i\vec x_i\right)
\le\sum_{i=1}^m\lambda_i f(\vec x_i).\]
\end{theoreme}
\begin{resultat}{Une extension de la définition}
Deux points deviennent une famille finie ; la moyenne des valeurs majore la valeur à la moyenne.
\end{resultat}
\vfill Les poids doivent être positifs ou nuls et leur somme doit valoir $1$.
\note{Cours original, pages PDF 114–115. Résultats admis ; interprétation sans démonstration.}
\end{frame}
